% !TeX root = proyecto.tex

%========================================================================================
% CAPÍTULO 10: TRABAJO FUTURO
%========================================================================================

\chapter{Trabajo Futuro}
\label{cap:trabajofuturo}

A pesar de la madurez tecnológica alcanzada por la arquitectura neuro-simbólica y la culminación exitosa de los prototipos delineados en el Trabajo Terminal II, la magnitud y urgencia del problema social abordado exigen una visión de escalabilidad continua. Este capítulo traza la hoja de ruta técnica a mediano y largo plazo, proponiendo las optimizaciones e integraciones de \textit{software} necesarias para transicionar el sistema de un entorno de validación hacia una infraestructura operativa a nivel nacional.

\section{Extracción de Entidades Nombradas (NER) bajo Principios Éticos}
\label{sec:tf_ner}

El siguiente hito evolutivo en el procesamiento de lenguaje natural del sistema consiste en la implementación de un módulo robusto de Extracción de Entidades Nombradas (NER, por sus siglas en inglés). Si bien BETO y BERTopic otorgan comprensión y contexto a nivel de documento, es imperativo extraer la granularidad de los actores involucrados. Se propone realizar un \textit{fine-tuning} sobre arquitecturas especializadas como \textit{spaCy} o integrar Modelos de Lenguaje de Gran Escala (LLMs) para minar entidades críticas como nombres de municipios, autoridades judiciales a cargo (jueces, fiscales) e instituciones gubernamentales responsables (e.g., el Desarrollo Integral de la Familia, DIF). 

Sin embargo, el rasgo distintivo de esta futura implementación será su diseño \textit{ético por defecto}. El módulo NER deberá ser programado con heurísticas de enmascaramiento estricto para ignorar, ofuscar o destruir deliberadamente cualquier token clasificado como nombre propio de un menor de edad. De este modo, se garantiza la extracción de inteligencia táctica para el seguimiento institucional, blindando al mismo tiempo la identidad y la integridad de las víctimas indirectas contra cualquier exposición accidental en las bases de datos.

\section{Generación Automática de Reportes (Exportación)}
\label{sec:tf_reportes}

El tablero interactivo actual posee un alto valor para el análisis en tiempo real, pero el trabajo de defensa de derechos humanos requiere materialización documental. El desarrollo futuro debe contemplar la creación de un subsistema de renderizado para la generación de expedientes descargables en formatos estandarizados, como PDF y hojas de cálculo (Excel). 

Esta capacidad de exportación automatizada dotará a las organizaciones civiles de una herramienta formidable. Podrán transformar instantáneamente los clústeres semánticos de BERTopic y las métricas de BETO en expedientes impresos, proveyendo evidencia estadística sólida y curada que puede ser directamente presentada ante comisiones legislativas, mesas de trabajo interinstitucionales o tribunales, eliminando las horas de transcripción manual que sufren los trabajadores sociales.

\section{API Pública e Interoperabilidad para Organizaciones Civiles}
\label{sec:tf_api}

La verdadera democratización de la tecnología radica en su interoperabilidad. Actualmente, el ecosistema funciona como un monolito cerrado al exterior, donde el usuario debe ingresar forzosamente al \textit{Dashboard} visual. Como trabajo a mediano plazo, se proyecta el desarrollo de una Interfaz de Programación de Aplicaciones (API) pública, implementada bajo el estándar RESTful o a través de GraphQL, protegida mediante autenticación robusta (OAuth 2.0).

La apertura de esta API permitirá que fundaciones como \textit{Futuro con Derechos} conecten sus propios sistemas internos, aplicaciones móviles o bases de datos directamente a nuestro núcleo neuro-simbólico. En lugar de adaptar sus flujos de trabajo a nuestra plataforma, podrán consumir bajo demanda los datos deduplicados, clasificados y libres de ruido, integrando la inteligencia de NNA directamente en los tableros propietarios que ya dominan y operan diariamente.

\section{Despliegue en la Nube (Cloud Architecture) y Escalabilidad}
\label{sec:tf_cloud}

El entorno de ejecución presente depende de la contenerización local en \textit{Docker}, lo cual restringe el poder de cómputo a los fierros del servidor anfitrión. Para sostener un monitoreo ininterrumpido sobre cientos de fuentes de todo el país, la arquitectura debe migrar obligatoriamente hacia una infraestructura de nube pública, como Amazon Web Services (AWS) o Google Cloud Platform (GCP).

Esta transición arquitectónica habilitará el uso de bases de datos relacionales completamente gestionadas y replicadas, como Amazon RDS para PostgreSQL, garantizando tolerancia a fallos y copias de seguridad continuas. Asimismo, la orquestación del núcleo de inteligencia artificial se vería profundamente beneficiada. Desplegar los pesados modelos de tensores de BETO y BERTopic bajo servicios de elasticidad computacional permitiría asignar Unidades de Procesamiento Gráfico (GPUs) bajo demanda, escalando la capacidad del análisis semántico automáticamente durante picos de contingencia noticiosa, y apagando los recursos cuando el flujo informativo disminuya, optimizando dramáticamente los costos operativos.

\section{Validación Piloto en Campo y Refinamiento UX/UI}
\label{sec:tf_piloto}

Finalmente, la culminación de cualquier herramienta de \textit{software} con impacto social no ocurre en el repositorio de código, sino en las manos de sus usuarios. El trabajo futuro ineludible consiste en sacar el sistema del entorno académico y someterlo a pruebas piloto exhaustivas directamente en campo. 

El despliegue en fase \textit{beta} con los trabajadores sociales, psicólogos y analistas de datos de las Organizaciones No Gubernamentales resulta indispensable. Solo a través del monitoreo de su interacción diaria se podrán capturar las métricas cualitativas necesarias para afinar la Experiencia de Usuario (UX) y la Interfaz de Usuario (UI). Comprender cómo leen los clústeres, qué información necesitan de un vistazo y qué botones resultan confusos, es el paso final para transformar este modelo de ingeniería en un aliado transparente y natural en la lucha diaria por la protección de las infancias en México.
